\chapter{Using Lists in Python}

\section{Introduction: Real-Life Hook}

Students often keep track of multiple items:  

\begin{itemize}
    \item Favorite snacks  
    \item Subjects in school  
    \item Friends’ names  
    \item Scores in multiple tests  
\end{itemize}

In Python, \textbf{lists} let us store multiple items in a single variable, making programs organized and efficient.

\section{Concept: What is a List?}

\begin{definition}{List}
A \textbf{list} in Python is an ordered collection of items, which can be numbers, strings, or even other lists.
\end{definition}

\subsection{Creating a List}

\begin{verbatim}
fruits = ["apple", "banana", "mango"]
marks = [80, 70, 95]
\end{verbatim}

- Lists are enclosed in square brackets \texttt{[]}  
- Items are separated by commas  
- Order matters; the first item has index 0

\section{Accessing Items in a List}

\begin{verbatim}
print(fruits[0])  # apple
print(marks[2])   # 95
\end{verbatim}

\subsection{Updating Items in a List}

\begin{verbatim}
fruits[1] = "orange"
print(fruits)  # ["apple", "orange", "mango"]
\end{verbatim}

\subsection{Adding Items to a List}

\begin{verbatim}
fruits.append("grapes")
print(fruits)  # ["apple", "orange", "mango", "grapes"]
\end{verbatim}

\subsection{Removing Items from a List}

\begin{verbatim}
fruits.remove("apple")
print(fruits)  # ["orange", "mango", "grapes"]
\end{verbatim}

\section{Looping Through a List}

\begin{verbatim}
for fruit in fruits:
    print("I like", fruit)
\end{verbatim}

Output:

\begin{verbatim}
I like orange
I like mango
I like grapes
\end{verbatim}

\section{Step-by-Step Diagram: List Operations}

\begin{center}
\begin{tikzpicture}[node distance=1.5cm]
\node (start) [draw, ellipse, fill=green!20] {Start Program};
\node (create) [draw, rectangle, fill=yellow!20, below of=start] {Create List};
\node (access) [draw, rectangle, fill=orange!20, below of=create] {Access / Update / Add / Remove Items};
\node (loop) [draw, rectangle, fill=orange!20, below of=access] {Loop Through List};
\node (output) [draw, ellipse, fill=green!20, below of=loop] {End Program};
\draw[->] (start) -- (create);
\draw[->] (create) -- (access);
\draw[->] (access) -- (loop);
\draw[->] (loop) -- (output);
\end{tikzpicture}
\end{center}

\section{Mini-Project: Favorite Foods List}

Create a program that stores your favorite foods in a list and prints them nicely:

\begin{verbatim}
foods = ["pizza", "burger", "ice cream"]
foods.append("chocolate")
for food in foods:
    print("I love", food)
\end{verbatim}

Output:

\begin{verbatim}
I love pizza
I love burger
I love ice cream
I love chocolate
\end{verbatim}

\section{Using Lists with User Input}

\begin{verbatim}
numbers = []
for i in range(3):
    num = int(input("Enter a number: "))
    numbers.append(num)
print("Your numbers are:", numbers)
\end{verbatim}

- User inputs are stored dynamically in the list  
- Loops make collecting multiple inputs easy

\section{Key Concepts}

\begin{itemize}
    \item Lists store multiple items in a single variable  
    \item Items are ordered and indexed starting from 0  
    \item Lists can be updated, added to, or removed from  
    \item Loops work perfectly with lists to process each item  
\end{itemize}

\section{End-of-Chapter Exercises}

\begin{exercise}{Very Short Questions}
\begin{enumerate}
    \item What is a list in Python?  
    \item How do you access the first item of a list?  
    \item Name two ways to modify a list.  
\end{enumerate}
\end{exercise}

\begin{exercise}{Short Answer Questions}
\begin{enumerate}
    \item Write a Python program to store 5 of your favorite movies in a list and print them.  
    \item How can you remove an item from a list? Give an example.  
    \item Explain why loops are useful with lists.  
\end{enumerate}
\end{exercise}

\begin{exercise}{Activity-Based Questions}
\begin{enumerate}
    \item Ask the user for 5 city names and store them in a list. Print them with a message "I want to visit \underline{\hspace{2cm}}".
    \item Create a list of test scores and calculate the average using a loop.  
    \item Make a list of your favorite subjects, update one of them, and print the final list.  
\end{enumerate}
\end{exercise}

\section{Teacher Notes}

\marginnote{
\textbf{Learning Objectives:}  
\begin{itemize}
    \item Students learn to create and manage lists.  
    \item Students can loop through lists for processing.  
    \item Students can update, add, and remove items dynamically.  
\end{itemize}

\textbf{Teaching Tips:}  
\begin{itemize}
    \item Use real-life examples (foods, movies, friends’ names).  
    \item Demonstrate append, remove, and looping live in Python IDE.  
\end{itemize}

\textbf{Assessment:}  
\begin{itemize}
    \item Check correctness of list operations.  
    \item Observe understanding of loops with lists.  
\end{itemize}
}